\section{Introduction}
\label{sec:intro}

Coding agents navigate repositories, read and edit files, and invoke builds
and tests through repeated tool-use loops~\cite{sweagent2024,openhands2025}.
Several tasks may begin with the same repository and dependencies while
requiring independent subsequent changes. Other tasks need only an ordinary
writable volume, or read-only access to a published environment. A file system
serving these uses must distinguish the semantics of a task's view from the
representation of its metadata and the geometry of its stored data.

These choices address different costs. Copying a baseline for every task
duplicates provisioning work and private state. Sharing that baseline leaves
open how clients obtain its directory entries, attributes, and data locations.
Storing those records in objects introduces a request-versus-transfer tradeoff.
The data path has a related but separate tradeoff: larger objects can amortize
requests, while smaller independently readable units can avoid fetching
unused bytes. A short application read may still cause a much larger object
transfer or physical storage read.

\sys is an object-backed file system that exposes three optional mechanisms
at these boundaries. \emph{Delta-based multi-agent workspaces} compose a
shared immutable baseline with private namespace and data-mapping state.
\emph{Packed Metadata} organizes immutable records into indexed, bounded
groups on object storage. \emph{Variable-Granularity Object Storage} controls
client object formation and, on a compatible backend, permits optional
encoding and verification layouts. Each mechanism has a useful role without
the other two: a workspace can inherit a native metadata layer, a packed
snapshot can be mounted read-only, and an ordinary volume can choose its
object block size while using the backend's default layout.

The mechanisms also compose. Many workspaces can reference one packed
baseline, amortizing its construction and reusing fetched pages. Their private
changes leave immutable baseline objects available to other readers. A packed
container or a newly uploaded data object may additionally use a compatible
backend layout. These interactions do not imply additive or multiplicative
gains. Frequent publication can erase packing's amortization, unrelated point
lookups can waste page bytes, and whole-object verification can prevent a
small file read from becoming a small backend range. Selection should follow
the workload and persistence contract.

The individual ideas have substantial precedents. OverlayFS provides layered
namespace semantics~\cite{overlayfs}, BranchFS uses copy-on-write branches for
agentic exploration~\cite{branchfs2026}, and DeltaFS separates private namespace
changes over shared storage~\cite{deltafs2021}. Immutable distribution systems
already support indexed, demand-loaded baselines~\cite{cvmfs2011,nydus_rafs}.
Our research question is therefore: \emph{where should view composition,
immutable metadata access, and object granularity be selected, and under
which workloads do their individual or combined costs justify those choices?}
Full directory copying is one control, not the only relevant baseline.

This paper develops the state and access rules behind these choices. It
separates a workspace's durable private state from transient dirty buffers,
and a baseline's logical identity from its metadata format. It also separates
the client object block size $g$, finalized object length $L_o$, backend
erasure-coding block size $B$, and physical verification/read quantum $q$.
The current implementation supports concrete subsets through different mount
paths; architectural optionality does not mean that every combination is
available as interchangeable runtime flags.

The evidence consists of an implementation audit and a reanalysis of
1,138,884 FUSE requests from eight coding-agent task runs. These observations
characterize request structure, rather than \sysname's performance. We specify
controlled experiments for each mechanism, legal combinations, lifecycle
costs, and complete tasks. The resulting contribution is a design and research
protocol whose hypotheses can be tested without confusing a change of runtime,
cache policy, or durability boundary with a mechanism's effect.
